\section{Methods}\label{sec:methodology}

\subsection{Analytical design and estimands}

The primary estimand is out-of-sample discrimination of four
high-confidence global LCZ weak labels from local climate-morphology features.
The analysis distinguishes transferable predictive signal from agreement with
local audited subtypes. Secondary estimands are (i) the
incremental predictive contribution of feature families, (ii) the paired
contribution of network 2SFCA thresholds, (iii) leave-district transfer, and
(iv) descriptive association between the weak classes, morphology, and summer
LST. Audited LCZ-like local subtype prediction is
assessed descriptively on the completed audit, but remains limited by the
small high-quality subset and the absence of a trained deep model.

The workflow is summarized in Figure~\ref{fig:workflow}. Data transformations,
model settings, and reported results are archived with the public analysis
package. Neural components are outside the present analysis and are clearly separated from the completed
tabular analysis.

\begin{figure}[!htbp]
    \centering
    \includegraphics[width=\textwidth]{fig_2_methodology_workflow.png}
    \caption{Completed evidence workflow and confirmatory status gate. Network
    accessibility is represented by network-cost 2SFCA, whereas the legacy
    Euclidean road-context fields are supplementary only.}
    \label{fig:workflow}
\end{figure}

\subsection{Grid construction and feature engineering}

The primary grid unit is 250 m. Raster variables were summarized using valid
pixels only; vector variables were aggregated by exact intersection where
semantically appropriate. Predictor fields were screened to remove identifiers,
row/column indices, district names, weak-label quality variables, LST response
variables, and valid-pixel counts. Missingness was interpreted before
imputation: building and floor variables are structurally absent in units
without buildings, and orientation or endpoint metrics are structurally absent
without road support. The modelling pipeline uses fold-specific median
imputation rather than global imputation.

The v8 matrix contains 374 fields, of which 294 were candidate numeric
predictors after target and quality-control exclusions. At absolute correlation
$r\geq0.95$, 490 pairs formed 40 connected groups. The first baselines therefore
use representative variables and family-level ablations rather than all
available fields. Table~\ref{tab:feature_families} lists the eight interpretive
families. LST is response-only; OSM semantic context is a proxy; height,
frontage, and endpoint fields remain explicitly qualified.

\begin{table*}[htbp]
\centering
\footnotesize
\caption{Feature families used in the evidence-grounded baseline.}
\label{tab:feature_families}
\resizebox{\textwidth}{!}{%
\begin{tabular}{lll}
\toprule
Family & Representative variables & Primary role \\
\midrule
2D morphology & coverage, density, compactness, fragmentation & urban fabric \\
3D/surface model & surface-model SD/range, floor-height proxies & surface roughness/vertical intensity \\
Vegetation & canopy cover, p95 height, volume, NDVI & 3D green structure \\
Green-blue/coast & NDWI, water/tree shares, coast distance & climate context \\
Street network & road/intersection/dead-end density, entropy & connectivity proxy \\
Street interface & continuity, gap, open-buffer share & frontage proxy \\
Green access & network-cost 2SFCA at 400/800/1200 m & competition-adjusted access \\
Semantic context & Dynamic World, WorldCover, OSM shares & proxy land-use context \\
\bottomrule
\end{tabular}%
}
\begin{minipage}{0.96\textwidth}\vspace{2pt}\scriptsize Note: LST and label-quality fields were excluded from predictors. OSM, floor count, frontage, and endpoint topology remain explicitly named proxies.\end{minipage}
\end{table*}



\subsection{Network-cost green-space 2SFCA}

Green-space context is defined by supply, demand, and shortest-path cost rather
than by a circular buffer. OSM park and garden polygon area provides supply
$S_j$. The 2024 neighbourhood population is allocated to 250 m units as demand
$P_i$ in proportion to residential floor-area proxy. Municipal road
centrelines are planarized after excluding explicitly coded motorways. Origins
and green-space access points are snapped to this network; their snap distances
are included in generalized distance $d_{ij}$.

For threshold $d_0$, step 1 calculates the supply--demand ratio

\begin{equation}
R_j = \frac{S_j}{\sum_{i:d_{ij}\leq d_0} P_i},
\end{equation}

and step 2 assigns each unit

\begin{equation}
A_i = \sum_{j:d_{ij}\leq d_0} R_j.
\end{equation}

The preregistered primary threshold is 800 m, with 400 and 1,200 m sensitivity
thresholds. Raw square metres per person are retained for description;
$\log(1+A_i)$ enters the model because of the long upper tail. This is an OSM
and municipal-road-network accessibility proxy, not a pedestrian routing
product, official public-green inventory, or measure of park quality.

To audit supply geography, the 2023-modified IZBB north/south park-maintenance
inventory was normalized by district and neighbourhood. It has no coordinates
or public-access field and is limited to branch responsibility. We therefore
compare log counts and log areas by district and use leave-one-district-out
correlation ranges. Name matching is used only as a quality check.

\subsection{Weak-label definition and audit sampling}

For unit $i$, the weak label is

\begin{equation}
l_i^{weak}=\arg\max_c A_{ic},
\end{equation}

where $A_{ic}$ is class-$c$ area or probability support. Weak confidence is

\begin{equation}
conf_i^{weak}=\max_c(A_{ic}/A_i).
\end{equation}

The manual-audit sample was allocated across weak class and preliminary
built-intensity strata. If $N_h$ is the size of stratum $h$, $N$ is the frame,
$n$ is the target sample, and $n_{min}$ protects rare strata, the initial
allocation follows

\begin{equation}
n_h=\max\{n_{min},\mathrm{round}(nN_h/N)\},
\end{equation}

with certainty selection where a stratum is smaller than its target. Spatial
balance prevents the sample from collapsing into common central fabrics. The
priority queue changes review order only; it does not change the sampling
frame, inclusion weights, or validation role.

\subsection{Baseline models and additive ablations}

Four tabular learners were evaluated as complementary baselines. Random Forest
used 300 trees, minimum leaf size 2, square-root feature sampling, and
balanced-subsample class weights. Extra Trees used the same tree count and leaf
size with balanced class weights. XGBoost used 300 histogram-based trees,
maximum depth 6, and learning rate 0.1 with fold-specific balanced sample
weights. LightGBM used 300 trees, learning rate 0.1, balanced class weights,
and gain-based feature importance. A majority classifier provided a floor.
Random seed 42 controls fold construction; fold-specific seeds control models.

Four nested recipes isolate information families. Recipe A contains 15 2D/3D
morphology variables. Recipe B adds seven vegetation-structure variables.
Recipe C adds seven coastal and green-blue variables. Recipe D adds remote
spectral and semantic variables, OSM proxy context, exact street-network
metrics, street-interface proxies, and one 2SFCA field, for 68 predictors. The
2SFCA experiment compares recipe D with the accessibility field omitted and
with otherwise identical 400, 800, or 1,200 m fields. Because semantic remote
sensing may overlap with the global LCZ weak-label production process, the
nested results are interpreted as discrimination evidence, not independent
causal contribution.

\subsection{Spatial validation and transfer testing}

The main validation uses five-fold StratifiedGroupKFold. Spatial groups are
5-by-5-grid-unit blocks, nominally 1.25 km, and no block appears in both training and
test data within a fold. Median imputation is fit only to each training fold.
Metrics are accuracy, balanced accuracy, weighted F1, macro-F1, and per-class
precision, recall, and F1. Out-of-fold probabilities provide the uncertainty
map.

A separate leave-one-district-out (LODO) transfer test evaluates the LightGBM recipe
D, with Random Forest as a comparison baseline. Districts with at least 100
eligible observations were held out one at a time, yielding 11 tests and 5,279
test predictions. Because several held-out
districts do not contain all four classes, LODO macro-F1 is calculated over
classes present in that district and is reported alongside balanced accuracy.
This is a harder transfer test, not a replacement for spatial blocks.

\subsection{Ablation, explainability, and uncertainty}

Additive ablations are compared on the same folds. For each 2SFCA threshold,
fold-level macro-F1 and balanced-accuracy differences are paired with the same
model without 2SFCA. Gain-based LightGBM importance is averaged across the five
primary fits. This importance is useful for screening but can be biased toward
correlated or high-cardinality predictors; it is not interpreted as a causal
effect. Out-of-fold maximum probability is mapped as a first uncertainty
indicator. A fuller multimodal study would additionally require calibrated
probabilities, SHAP for tabular components, and spatial attribution for neural
branches.

\subsection{Climate-response validation}

LST is never a morphology predictor. For the 5,339-unit analysis subset, we
report class-specific medians and interquartile ranges, a Kruskal--Wallis test,
and epsilon-squared effect size

\begin{equation}
\epsilon^2=\frac{H-k+1}{n-k},
\end{equation}

where $H$ is the Kruskal statistic, $k$ is the number of classes, and $n$ is the
valid sample. Spearman correlations connect selected morphology, vegetation,
network, coast, and 2SFCA variables to LST. These are bivariate descriptive
associations. They do not estimate cooling causality or air-temperature
exposure.

\subsection{Confirmatory CLIMORFA architecture and stopping rule}

The confirmatory architecture has three branches: (i) raster chips for
Sentinel, land-cover probabilities, and the 5 m surface model; (ii) multichannel
figure-ground chips for buildings, streets, open space, vegetation, and height
coding; and (iii) tabular morphometric, vegetation, coast, network, and 2SFCA
features. Each branch creates an embedding $e_b$. Late fusion concatenates
embeddings,

\begin{equation}
z_i=[e_{r,i};e_{g,i};e_{t,i}],
\end{equation}

or uses learned gates $\alpha_b$,

\begin{equation}
z_i=\sum_{b\in\{r,g,t\}}\alpha_b W_b e_{b,i}.
\end{equation}

Mandatory ablations remove the surface model, vegetation, network, figure-ground, and remote
semantic inputs. Integrated Gradients, Grad-CAM or occlusion, and tabular SHAP
are required. These analyses cannot enter the
abstract, headline tables, or conclusion as a deep-model performance claim
until the branches are trained and evaluated on an audited spatial-validation
design. This paper therefore reports an evidence-grounded tabular
baseline and bounded audit validation, not CLIMORFA deep-model performance.
